% =====================================================================
\appendix
\addto\captionsspanish{\renewcommand{\appendixname}{Anexo}}
\renewcommand{\appendixname}{Anexo}
\chapter{Reproducibilidad y organización del proyecto}
\label{anx:repro}

El proyecto está organizado según las fases de CRISP-DM: cada fase es una carpeta, y adentro hay una
subcarpeta por tema (tabla~\ref{tab:carpetas}). Cada subcarpeta tiene siempre lo mismo: un
\texttt{README.md} que dice qué hace, qué recibe y qué entrega; su script; una carpeta
\texttt{salidas/} con los resultados; y un \texttt{metricas.csv} con los números que usan este
informe y la presentación. Ningún número del informe se escribió a mano: el script
\texttt{numeros.py} los toma de esos archivos al compilar.

{\footnotesize
\begin{longtable}{@{}L{5.4cm}L{8.4cm}@{}}
\toprule
\textbf{Carpeta} & \textbf{Qué contiene} \\
\midrule
\endhead
\multicolumn{2}{@{}l}{\textit{1\_comprension\_del\_problema}} \\
\carpeta{hipotesis_y_objetivos} & Hipótesis, criterio de aceptación y objetivos \\
\midrule
\multicolumn{2}{@{}l}{\textit{2\_comprension\_de\_los\_datos}} \\
\carpeta{noticias} & Corpus, scripts de recolección y auditoría \\
\carpeta{mercado} & Precios de las acciones y selección del panel \\
\midrule
\multicolumn{2}{@{}l}{\textit{3\_preparacion\_de\_los\_datos}} \\
\carpeta{regla_17} & Cada nota a su rueda \\
\carpeta{sentimiento_finbeto} & Etiquetado de Gemma, FinBETO de tres clases y clasificación del corpus \\
\carpeta{sentimiento_por_rueda} & Sentimiento de cada rueda y sentimiento inesperado \\
\carpeta{sentimiento_por_empresa} & Cruce de entidades con el sentimiento y sentimiento de cada empresa \\
\carpeta{vix_argentino} & VIX argentino y su escala \\
\carpeta{regimen_hmm} & Regímenes del mercado (solo para el anexo) \\
\carpeta{dataset_unificado} & Las tablas que usan los modelos y el diccionario de variables \\
\midrule
\multicolumn{2}{@{}l}{\textit{4\_modelado}} \\
\carpeta{noticias_a_precio} & Modelo A contra B: ¿las noticias anticipan al precio? \\
\carpeta{precio_a_sentimiento} & Modelo A contra B: ¿el precio anticipa al sentimiento? \\
\midrule
\multicolumn{2}{@{}l}{\textit{5\_evaluacion}} \\
\carpeta{veredicto_hipotesis} & El criterio de aceptación aplicado a los resultados \\
\carpeta{anexo} & Robustez y verificación con VAR y Granger \\
\midrule
\multicolumn{2}{@{}l}{\textit{6\_despliegue}} \\
\carpeta{informe} & Este documento \\
\carpeta{presentacion} & Guion de la presentación \\
\bottomrule
\caption{Organización del proyecto. Fuente: elaboración propia.}
\label{tab:carpetas}
\end{longtable}
}

Cada paso de la preparación verifica, además, que reproduce los resultados de la versión anterior
cuando se le dan sus mismas entradas. Por ejemplo, con las etiquetas de cuatro clases, el
sentimiento por rueda sale idéntico al anterior. Así se sabe que lo único que cambió es lo que se
decidió cambiar.

% =====================================================================
\chapter{Hiperparámetros elegidos}
\label{anx:hiper}

La tabla~\ref{tab:hiper} muestra los hiperparámetros que eligió la validación cruzada temporal dentro
del entrenamiento (sección~\ref{sec:diseno}), y el AUC promedio de esa validación.

\begin{table}[H]
\centering\small
\begin{tabular}{@{}llccc@{}}
\toprule
 & & \textbf{Logística} & \textbf{Random forest} & \textbf{AUC de validación} \\
\textbf{Dirección} & \textbf{Modelo} & \textbf{$C$} & \textbf{profundidad / hoja} & \textbf{(logística)} \\
\midrule
Noticias \flecha{} precio & A & \dato{np.h.logistica.A.C} & \dato{np.h.random_forest.A.max_depth} / \dato{np.h.random_forest.A.min_samples_leaf} & \dato{np.h.logistica.A.auc_validacion@3} \\
 & B & \dato{np.h.logistica.B.C} & \dato{np.h.random_forest.B.max_depth} / \dato{np.h.random_forest.B.min_samples_leaf} & \dato{np.h.logistica.B.auc_validacion@3} \\
Precio \flecha{} sentimiento & A & \dato{ps.h.logistica.A.C} & \dato{ps.h.random_forest.A.max_depth} / \dato{ps.h.random_forest.A.min_samples_leaf} & \dato{ps.h.logistica.A.auc_validacion@3} \\
 & B & \dato{ps.h.logistica.B.C} & \dato{ps.h.random_forest.B.max_depth} / \dato{ps.h.random_forest.B.min_samples_leaf} & \dato{ps.h.logistica.B.auc_validacion@3} \\
\bottomrule
\end{tabular}
\caption{Hiperparámetros elegidos. Fuente: elaboración propia (\carpeta{4_modelado}).}
\label{tab:hiper}
\end{table}

% =====================================================================
\chapter{Los regímenes de mercado}
\label{anx:hmm}

El lado bursátil clasificó cada rueda en un régimen de mercado con un modelo oculto de Markov: un
modelo que supone que el mercado pasa por unos pocos estados que no se observan directamente, con
una media y una varianza del VIX propias, y que salta de uno a otro según una matriz de
probabilidades de transición \citep{hamilton1989,kim1994,rabiner1989}. Se estimó con
\texttt{MarkovRegression} de \texttt{statsmodels} sobre las \dato{hmm.ruedas_total} ruedas del VIX
argentino. Como el régimen se estima con el mismo VIX que ya está en el modelo, en esta versión no
entra en el modelo principal: se usa solo como prueba de robustez (sección~\ref{sec:robustez}), con
la versión \textit{filtrada}, que solo usa información hasta cada día \citep{majumder2023}. Coincide
con la suavizada en el \dato{hmm.coincidencia_filtrado_pct}\% de los días.

El número de regímenes se eligió con el criterio bayesiano de información \citep{schwarz1978},
descartando los ajustes que no convergen o que crean un régimen de menos de cinco días
(tabla~\ref{tab:bic}). Con tres regímenes el BIC es el menor entre los ajustes válidos
(tablas~\ref{tab:regimenes} y~\ref{tab:transicion}). Las \dato{np.ruedas_test} ruedas del período
de prueba están todas en el régimen Calmo.

\begin{table}[H]
\centering\small
\begin{tabular}{@{}ccccc@{}}
\toprule
$k$ & \textbf{BIC} & \textbf{Convergió} & \textbf{Régimen más chico} & \textbf{Válido} \\
\midrule
\tabla{bic}
\bottomrule
\end{tabular}
\caption{Selección del número de regímenes. Fuente: elaboración propia
(\carpeta{3_preparacion_de_los_datos/regimen_hmm}).}
\label{tab:bic}
\end{table}

\begin{table}[H]
\centering\small
\begin{tabular}{@{}lrrrr@{}}
\toprule
\textbf{Régimen} & \textbf{VIX promedio} & \textbf{Ruedas (total)} & \textbf{Ruedas de la ventana} & \textbf{Ruedas (filtrado)} \\
\midrule
\tabla{regimenes}
\bottomrule
\end{tabular}
\caption{Los tres regímenes. Fuente: elaboración propia
(\carpeta{3_preparacion_de_los_datos/regimen_hmm}).}
\label{tab:regimenes}
\end{table}

\begin{table}[H]
\centering\small
\begin{tabular}{@{}lccc@{}}
\toprule
\textbf{De \flecha{} a} & \textbf{Calmo} & \textbf{Moderado} & \textbf{Tenso} \\
\midrule
\tabla{transicion}
\bottomrule
\end{tabular}
\caption{Probabilidad de pasar de un régimen (fila) a otro (columna) de una rueda a la siguiente.
Fuente: elaboración propia (\carpeta{3_preparacion_de_los_datos/regimen_hmm}).}
\label{tab:transicion}
\end{table}
